% !PRINTED 380
\weberpage{380}
%<<CONT
rationally determined by these $a_{k,i}$, and from (16) and (17) one then obtains $\xi_{4i}$, $\xi_{5i}$, $\xi_{6i}$, $\xi_{7i}$, $x_4$, $x_5$, $x_6$, $x_7$.

By substituting the $\xi_1$, $\xi_2$, $\xi_3$ into equation (4) one obtains the equation of a curve of the fourth order, whose coefficients are rational functions of the $a_{k,i}$, and the discriminant of this equation cannot vanish identically, because conversely, as we have seen, one can derive a system of equations (18), (19), (20) from the equation of a curve of the fourth order without a singular point.

From the equations (18), (16) there then also follow the formulae (7), and those which arise from them by permutation of 4 with 5, 6, 7; from which it is to be concluded that $x_4$, $x_5$, $x_6$, $x_7$ together with $x_1$, $x_2$, $x_3$ form a complete seven-system.

\S.\,101.

\gesperrt{The Galois group of the double tangent problem.}

The theorems that we have derived are sufficient to determine the \gesperrt{Galois} group of the algebraic equation on which the double tangents depend. We here take as \gesperrt{domain of rationality the field which consists of all rational functions of the 14 ratios of the coefficients of a general ternary form of degree $4$}, in which these coefficients count as \gesperrt{independent variables}. The equation of degree $28$ we may then think of as formed, say, in the following way: we take as unknowns the abscissae of the points of intersection of the double tangents with an arbitrary fixed straight line $L$. By means of the roots $\xi$ of this equation, which we call the \gesperrt{double tangent equation}, the double tangents are then rationally representable.

If one uses a \gesperrt{Cartesian} coordinate system $x$, $y$, whose $x$-axis is the line $L$, the equation of a double tangent takes the form
\weq{1}{y = \Theta\,(x - \xi),}
and the double tangent equation can be formed, if $F(x,y) = 0$ is the equation of the curve of the fourth order, by seeking the conditions that the function of $x$
%CONT>>
